\subsection{Quasienergy and Floquet States under $H_{\text{disp}}$}
In this subsection we relate an approximation of quasienergy and Floquet states to a rotating-wave approximation.
\subsubsection{Quasienergy Operator for the Driven Dispersive Hamiltonian}

The driven dispersive Hamiltonian is:
\begin{equation}
\begin{split}
\label{Driven dispersive}
    H(t) &= \bar{\omega}_b b^\dagger b + \frac{K}{2} b^{\dagger 2} b^2 + \bar{\omega}_c c^\dagger c \\
    &\quad + \chi b^\dagger b c^\dagger c + 2\Omega_0 \cos(\omega_d t) (b + b^\dagger).
\end{split}
\end{equation}
Following the recipe in the previous section, the quasienergy operator in the extended Hilbert space $\mathcal{F}$ is written as
\begin{equation}
    \overline{Q} = \overline{Q}_0 + \overline{V}.
\end{equation}
A convenient basis is $\kett{\mu,(n,m)}$, where $\kett{\mu,(n,m)}(t) = e^{i\mu\omega_dt}|n,m\rangle$ maps time to basis of the static dispersive Hamiltonian. The unperturbed component $\overline{Q}_0$ is diagonal in the basis $\kett{\mu,(n,m)}$:
\begin{widetext}
\begin{align}
\overline{Q}_0 &= \sum_{\mu,n,m} \left( \bar{\omega}_b n + \frac{K}{2} n(n-1) + \bar{\omega}_c m + \chi nm + \mu \omega_d \right)
\kett{\mu,(n,m)}\bbra{\mu,(n,m)}, \\
\overline{V} &= \Omega_0 \sum_{\mu,n,m} \sqrt{n+1}
\left( \kett{\mu+1,(n+1,m)}\bbra{\mu,(n,m)}
+ \kett{\mu-1,(n+1,m)}\bbra{\mu,(n,m)} \right)
+ \text{h.c.}
\end{align}
\end{widetext}

The perturbation $\overline{V}$ arises from the periodic drive,
\begin{equation}
    2\Omega_0 \cos(\omega_d t)=\Omega_0\!\left(e^{i\omega_d t}+e^{-i\omega_d t}\right),
\end{equation}
which mixes adjacent photon sectors. In the extended (Floquet) Hilbert space, the Fourier components $e^{\pm i\omega_d t}$ connect adjacent sectors $\mu\to\mu\mp 1$, while the drive excites FTT from $n\to n+1$. These processes have detunings
\begin{align}
    \Delta E_1 &= (\bar{\omega}_b - \omega_d)n + (n-1)K + \chi m,\\
    \Delta E_2 &=(\bar{\omega}_b + \omega_d)n + (n-1)K + \chi m,
\end{align}
respectively.
For our parameter regime of interest $|\bar{\omega}_b - \omega_d|\ll\omega_d$, we have parameters hierarchy
\begin{equation}
\label{hierarchy}
    \Omega_0 < \Delta E_1 \ll \Delta E_2.
\end{equation}
so that $\overline{V}$ can be treated perturbatively.

\subsubsection{Approximation of Quasienergy and Floquet States}
To show eigenvalue of $H_R$ can approximate quasienergies, we block-diagonalize the operator $\bar Q$  via a Schrieffer--Wolff transformation within the near-resonant subspace selected by the hierarchy in Eq.~\eqref{hierarchy}. Specifically, we define the projectors
\begin{equation}
    \bar P_\mu=\sum_{n,m}\ket{\mu-n,n,m}\bra{\mu-n,n,m},
\end{equation}
which collect the states coupled most strongly by the drive. Since $\bar P_\mu$ mixes different $\mu$ indices, it is convenient to relabel the basis with the unitary
\begin{equation}
    \bar R=\sum_{\mu,n,m}\ket{\mu,n,m}\bra{\mu-n,n,m}.
\end{equation}
In this relabeled frame,
\begin{align}
    \bar P_\mu^{R} &= \bar R^\dagger \bar P_\mu \bar R
    =\sum_{n,m}\ket{\mu,n,m}\bra{\mu,n,m},\\
    \bar Q_{R} &= \bar R^\dagger \bar Q\,\bar R,
\end{align}
so that $\bar Q_R$ acquires a block structure in the fixed-$\mu$ sectors:
\begin{widetext}
\begin{equation}
    \bar q_R=
    \begin{pmatrix}
     \ddots & \vdots & \vdots & \vdots & \vdots & \vdots & \ddots \\
     \cdots & [H_R-2\omega_d I] & 0 & [b] & 0 & 0 & \cdots \\
     \cdots & 0 & [H_R-\omega_d I] & 0 & [b] & 0 & \cdots \\
     \cdots & [b^\dagger] & 0 & [H_R] & 0 & [b] & \cdots \\
     \cdots & 0 & [b^\dagger] & 0 & [H_R+\omega_d I] & 0 & \cdots \\
     \cdots & 0 & 0 & [b^\dagger] & 0 & [H_R+2\omega_d I] & \cdots \\
     \ddots & \vdots & \vdots & \vdots & \vdots & \vdots & \ddots
    \end{pmatrix},
\end{equation}
\end{widetext}
where $[O]$ denotes the matrix of $O$ in the $\{\ket{n,m}\}$ basis. To first order, the effective quasienergy operator is given by the block-diagonal part,
\begin{equation}
    \sum_\mu \bar P_\mu^{R}\,\bar Q_R\,\bar P_\mu^{R}.
\end{equation}
It therefore suffices to diagonalize a single sector (e.g., $\mu=0$), denoted by $H_R$:
\begin{equation}
    H_R\ket{\psi_j}=E_j\ket{\psi_j},
    \qquad
    \ket{\psi_j}=\sum_{n,m}c_{nm}^j\ket{n,m}.
\end{equation}
The quasienergy  is then
\begin{equation}
    \varepsilon_{j,\mu}=E_j+\mu\omega_d,
\end{equation}
with corresponding Floquet states (in the relabeled frame)
\begin{equation}
    \ket{\psi_{j,\mu}}_{R}=\sum_{n,m}c_{nm}^j\ket{\mu,n,m}.
\end{equation}
\subsubsection{Connection with the RWA}
The relabeling $\bar R$ in extended Hilbert space implements the same rotation as a standard rotating-frame transformation in the physical Hilbert space. Acting on a Floquet basis state and evaluating its time dependence gives
\begin{align}
    \bigl(\bar R\ket{\mu,n,m}\bigr)(t)
    &= \ket{\mu+n,n,m}(t)
     = \ket{n,m}\,e^{i(\mu+n)\omega_d t} \nonumber\\
    &= \bigl(e^{in\omega_d t}\ket{n,m}\bigr)\,e^{i\mu\omega_d t},
\end{align}
so $\bar R$ amounts to multiplying $\ket{n,m}$ by the phase $e^{in\omega_d t}$. This is precisely the action of the rotating-frame unitary
\begin{equation}
    U_R(t)=e^{i\omega_d t\, b^\dagger b},
\end{equation}
which yields the effective Hamiltonian
\begin{equation}
    H_{\mathrm{eff}}(t)=U_R^\dagger(t)H(t)U_R(t)-i\,U_R^\dagger(t)\dot U_R(t).
\end{equation}
In this frame, terms oscillating at $\pm 2\omega_d$ are rapidly varying; neglecting them under the rotating-wave approximation produces a time-independent Hamiltonian $H_R$.

